\documentclass[10pt,a4paper]{article}
\usepackage[margin=1.7cm]{geometry}
\usepackage[T1]{fontenc}
\usepackage{lmodern}
\usepackage{booktabs}
\usepackage{enumitem}
\usepackage{listings}
\usepackage{xcolor}
\usepackage{titlesec}

\setlist{nosep,leftmargin=1.2em}
\setlength{\parindent}{0pt}
\setlength{\parskip}{3pt}
\titlespacing*{\section}{0pt}{6pt}{2pt}
\titlespacing*{\subsection}{0pt}{4pt}{1pt}
\titleformat*{\section}{\large\bfseries}
\titleformat*{\subsection}{\normalsize\bfseries}
\lstset{basicstyle=\ttfamily\footnotesize,columns=fixed,basewidth=0.5em,keepspaces=true,
  frame=single,framerule=0.3pt,rulecolor=\color{gray},
  aboveskip=3pt,belowskip=3pt,xleftmargin=2pt,xrightmargin=2pt}
\newcommand{\code}[1]{\texttt{#1}}

\begin{document}
\sloppy

\begin{center}
{\Large\bfseries CS3423 Compilers II -- Mini-Assignment 5: Extending Kaleidoscope}\\[2pt]
CS24BTECH11030 \quad CS24BTECH11036
\end{center}

\section{Starting point}
Our base is the Kaleidoscope chapter-3 front end (LLVM IR generation, no JIT),
ported to \textbf{flex/bison}: \code{lexer.l} replaces \code{gettok()},
\code{parser.y} replaces the recursive-descent parser and the
\code{BinopPrecedence} table (now \code{\%left} declarations), \code{ast.h}
holds the AST classes shared by both translation units, and \code{toy.cpp}
holds code generation and the driver. Build with \code{make}; run with
\code{./toy < tests.kal}. All changes below are relative to that unmodified
port (git commit \code{388e72b}); together they touch 4 source files
(+341/$-$88 lines), mostly \code{toy.cpp}.

\section{Question 1: division (\code{/}) and greater-than (\code{>})}
\begin{itemize}
\item \textbf{Lexer:} no change needed, because single characters fall through the
  catch-all rule \code{.~\{ return yytext[0]; \}}, as in the original \code{gettok()}.
\item \textbf{Parser:} new rules \code{expr '/' expr} and \code{expr '>' expr};
  precedence \code{\%left '<' '>'} (lowest) $<$ \code{\%left '+' '-'} $<$
  \code{\%left '*' '/'}, all left-associative. Bison reports no conflicts.
\item \textbf{Codegen:} \code{int}: \code{sdiv}, \code{icmp sgt}/\code{slt}
  (\emph{signed}, since \code{int} is signed); \code{double}: \code{fdiv},
  \code{fcmp ogt}/\code{olt}. We chose \emph{ordered} compares (false if either
  operand is NaN, as in C) instead of the tutorial's unordered \code{ult}.
\item \textbf{Design choice:} a comparison always yields \code{bool} (\code{i1}).
  The tutorial's \code{uitofp} back to \code{double} was removed.
\item \textbf{Semantic check:} an integer division whose divisor is the constant
  0 is rejected at compile time, using \code{dyn\_cast<ConstantInt>}.
\end{itemize}

\section{Question 2: \code{int} and \code{bool} types}
\subsection{Syntax (lexer + parser)}
\begin{itemize}
\item Keywords \code{int}, \code{double}, \code{bool} and literals \code{true},
  \code{false}. Flex prefers the longest match and then the earliest rule, so
  \code{int} is a keyword while \code{integer} remains an identifier.
\item Literals: \code{12} is an \code{int} (parsed with \code{strtol}; an
  out-of-range value is reported and the lexer returns \code{YYerror}, so the
  statement is discarded); \code{1.5}, \code{1.} and \code{.5} are \code{double}.
\item \textbf{Every variable definition is typed.} In this chapter, function
  parameters are the only variable definitions (there is no \code{var}/assignment),
  so each parameter carries a type. Functions also declare an \emph{explicit
  return type}:
\end{itemize}
\begin{lstlisting}
prototype ::= type id '(' (type id)* ')'    def int add(int a int b) a + b;
type      ::= int | double | bool            extern double sin(double x);
expr      ::= ... | type '(' expr ')'        def int trunc(double x) int(x);
\end{lstlisting}
The cast rule starts with a keyword token, so it cannot clash with calls
(\code{id '(' args ')'}) or parentheses, and bison still reports no conflicts.
Top-level expressions remain anonymous \code{double} functions; their value is
wrapped in an explicit \code{double(...)} cast, so \code{4 > 3;} yields \code{1.0}.

\subsection{Type representation and the symbol table}
\code{ast.h} gains \code{enum class VarType \{Double, Int, Bool\}}, a
\code{Param \{VarType Ty; std::string Name;\}} for prototype parameters (it
replaces the earlier stringly-typed \code{vector<vector<string>>}), a
\code{VarType RetTy} in \code{PrototypeAST}, a new \code{CastExprAST} node, and
\code{toLLVMType()}: \code{double}$\to$\code{double}, \code{int}$\to$\code{i32},
\code{bool}$\to$\code{i1}. The symbol table \code{NamedValues} changes from
\code{map<string, Value*>} to \code{map<string, SymbolInfo>}, where
\code{SymbolInfo = \{Value *Val; VarType Ty;\}}. It is filled from the
prototype's declared parameter types. Lookup uses \code{find()} instead of
\code{operator[]}, which used to silently insert unknown names.

\subsection{Type checking}
Expression types are read back from the generated LLVM value
(\code{isDoubleTy()}, \code{isIntegerTy(32)}, \code{isIntegerTy(1)}). The rules are:

\begin{center}\small
\begin{tabular}{@{}lll@{}}
\toprule
Construct & Allowed operand types & Result / action \\
\midrule
\code{+ - * /} & int,int \; | \; double,double \; | \; int,double (int promoted) & operand type \\
\code{< >} & same as above & \code{bool} \\
any binary op & either operand \code{bool} & error: ``operands cannot be bool'' \\
call \code{f(e$_1$,..)} & each $e_i$ must equal / implicitly convert to param $i$ & else error naming the argument \\
function body & must equal / implicitly convert to the declared return type & else error \\
\bottomrule
\end{tabular}
\end{center}

\subsection{Type casting}
A single helper, \code{castTo(V, Target, Explicit, Ctx)}, is used by binary
operators, call arguments, \code{return} values and \code{CastExprAST}. Only the
widening \code{int}$\to$\code{double} conversion is implicit. Everything else
needs the explicit \code{T(e)} syntax, so \code{bool} never leaks into
arithmetic by accident:

\begin{center}\small
\begin{tabular}{@{}llll@{}}
\toprule
From $\to$ To & Implicit? & Explicit & LLVM instruction \\
\midrule
int $\to$ double & yes & yes & \code{sitofp} \\
double $\to$ int & no (narrowing) & yes & \code{fptosi} (truncates toward 0) \\
bool $\to$ int / double & no & yes & \code{zext} / \code{uitofp} (true $\to$ 1, not $-1$) \\
int / double $\to$ bool & no & yes & \code{icmp ne 0} / \code{fcmp une 0.0} (NaN is true, as in C) \\
\bottomrule
\end{tabular}
\end{center}

\subsection{Other semantic checks added}
\begin{itemize}
\item Duplicate parameter names: \code{def int f(int a double a) a;} reports
  ``redefinition of parameter `a'\,''.
\item Extern/definition consistency: a definition (or another \code{extern})
  whose signature differs from an earlier declaration is rejected. Before this,
  LLVM silently created a renamed \code{g.1}. A definition whose body fails to
  type-check now keeps the earlier \code{extern} instead of erasing it.
\item \code{verifyFunction(*F, \&errs())} is now checked. Before this change,
  invalid IR such as \code{fadd i1} or a \code{double} function returning
  \code{i32} was printed silently. In our test suite the verifier never fires,
  because every bad case is caught by our own type checks first.
\item Error messages now name the variable, function, argument position and
  types involved.
\end{itemize}

\section{Summary of edits per file}
\begin{center}\small
\begin{tabular}{@{}lp{0.8\textwidth}@{}}
\toprule
File & Changes relative to the original port \\
\midrule
\code{lexer.l} & \code{int}/\code{double}/\code{bool}/\code{true}/\code{false} tokens; the single
  \code{NUMBER} token is split into \code{INTEGER} (\code{strtol}, range check) and
  \code{DOUBLE} (including the \code{.5} form). \\
\code{parser.y} & \code{\%union} gains \code{VarType ty}, \code{bool cond}, \code{int},
  \code{vector<Param>*}; new \code{type} rule; typed \code{params}; a return type
  on \code{prototype}; the \code{/}, \code{>} and cast rules; updated precedence. \\
\code{ast.h} & \code{VarType}, \code{Param}, \code{toLLVMType()}; typed literals in
  \code{NumberExprAST}; new \code{CastExprAST}; \code{RetTy} and getters in
  \code{PrototypeAST}. \\
\code{toy.cpp} & \code{SymbolInfo} symbol table; \code{castTo()}; typed
  \code{BinaryExprAST}/\code{CallExprAST}/return codegen; prototype signature and
  duplicate-parameter checks; checked \code{verifyFunction}; anonymous expressions
  cast to \code{double}. \\
\code{tests.kal} & New file: regression inputs for every rule in this report. \\
\bottomrule
\end{tabular}
\end{center}

\section{Testing}
\code{tests.kal} (in the directory) contains 32 cases covering every rule above.
Each case either produces valid IR or one of our own error messages. Examples:

\begin{lstlisting}
def double mix(int a double b) a + b;  ->  %casttmp = sitofp i32 %a to double
                                            %addtmp  = fadd double %casttmp, %b
def bool neg(int a) (a - 5) < 0;       ->  %cmptmp = icmp slt i32 %subtmp, 0 ; ret i1
def int bad1(bool a bool b) a + b;     ->  Error: operator '+': operands cannot be bool (got bool and bool)
def int bad2(double x) x;              ->  Error: return value of 'bad2': cannot implicitly convert double to int
idiv(7, 2.5);                          ->  Error: argument 2 of 'idiv': cannot implicitly convert double to int
extern int g(int x); def double g(double x) x;  ->  Error: conflicting types for 'g'
\end{lstlisting}

\textbf{Limitations.} There is no unary minus (write \code{0 - x}), comparing two
\code{bool}s is rejected rather than defined, and top-level expressions always
evaluate to \code{double}.

\end{document}
